% !TeX root = ../proyecto.tex

\begin{UseCase}{CU-02}{Gestión del Catálogo de Fuentes de Información}{%
	Permite al administrador gestionar el catálogo de fuentes periodísticas (RSS, HTML, sitemap) que alimentan el pipeline de recolección. Contempla el alta, edición, eliminación lógica, activación/desactivación y sondeo previo de URLs.
}
	\UCitem{Versión}{\color{Gray}
		2.0
	}\UCitem{Autor}{\color{Gray}
		Tesista de Ingeniería
	}\UCitem{Supervisa}{\color{Gray}
		Director de Trabajo Terminal
	}\UCitem{Actor}{
		Administrador del Sistema
	}\UCitem{Propósito}{\begin{Titemize}
		\Titem Mantener actualizado el catálogo de orígenes periodísticos accesibles por el recolector, garantizando cobertura geográfica y temática del corpus.
	\end{Titemize}
	}\UCitem{Entradas}{\begin{Titemize}
		\Titem URL de la fuente (obligatoria).
		\Titem Nombre del medio de comunicación (opcional; se infiere del dominio si se omite).
		\Titem Tipo de fuente (\texttt{auto} / \texttt{rss} / \texttt{sitemap} / \texttt{html}).
		\Titem Notas adicionales (campo libre, opcional).
	\end{Titemize}
	}\UCitem{Origen}{\begin{Titemize}
		\Titem Formulario y tabla de administración en la pantalla \IUref{IU-02}{Gestión de Fuentes}.
	\end{Titemize}
	}\UCitem{Salidas}{\begin{Titemize}
		\Titem Confirmación visual con los datos de la fuente creada, actualizada o eliminada.
		\Titem Resultado del sondeo técnico de la URL (tipo detectado, feeds RSS encontrados, crawl-delay, accesibilidad WAF).
		\Titem Tabla de fuentes actualizada con estado \texttt{ok / warning / blocked / error}.
	\end{Titemize}
	}\UCitem{Destino}{\begin{Titemize}
		\Titem Tabla \texttt{news\_sources} en PostgreSQL.
		\Titem Interfaz web (Pantalla \IUref{IU-02}{Gestión de Fuentes}).
	\end{Titemize}
	}\UCitem{Precondiciones}{\begin{Titemize}
		\Titem El actor debe contar con una sesión activa en el Dashboard web (CU-07).
		\Titem La base de datos PostgreSQL debe estar operativa.
	\end{Titemize}
	}\UCitem{Postcondiciones}{\begin{Titemize}
		\Titem Las fuentes activas quedan disponibles para el recolector en la próxima ejecución del pipeline (CU-01).
		\Titem Las fuentes predeterminadas del sistema solo pueden desactivarse; no pueden eliminarse.
	\end{Titemize}
	}\UCitem{Errores}{\begin{Titemize}
		\Titem {\bf \hypertarget{CU-02.E1}{E1}}: URL duplicada ya registrada en el catálogo \Trayref{CU-02}{A}.
		\Titem {\bf \hypertarget{CU-02.E2}{E2}}: Intento de eliminar una fuente predeterminada del sistema \Trayref{CU-02}{B}.
		\Titem {\bf \hypertarget{CU-02.E3}{E3}}: La URL no inicia con \texttt{http://} o \texttt{https://} \Trayref{CU-02}{C}.
	\end{Titemize}
	}\UCitem{Tipo}{
		Caso de uso primario
	}\UCitem{Observaciones}{
		El sondeo técnico (endpoint \texttt{POST /api/sources/probe}) es una operación no persistente: detecta el tipo de fuente y la accesibilidad WAF \textit{antes} de registrarla, permitiendo al administrador tomar una decisión informada.
	}
\end{UseCase}

%--------------------------------------
\begin{UCtrayectoria}
	\UCpaso[\UCActor] Ingresa a la pantalla de \IUref{IU-02}{Gestión de Fuentes} desde el menú lateral.
	\UCpaso[\UCActor] Introduce la URL de la nueva fuente en el formulario de alta.
	\UCpaso[\UCActor] Hace clic en \IUbutton{Sondear URL} para obtener una vista previa técnica antes de registrar.
	\UCpaso[\UCsist] Invoca \texttt{DynamicScraper.probe\_url()} que detecta el tipo de contenido (RSS, sitemap, HTML), verifica el archivo \texttt{robots.txt} y determina si requiere modo stealth \ErrorRef{CU-02}{E3}{URL inválida}.
	\UCpaso[\UCsist] Devuelve el resultado del sondeo al formulario (tipo detectado, feeds RSS encontrados, crawl-delay recomendado, accesibilidad WAF).
	\UCpaso[\UCActor] Revisa los resultados, ajusta el nombre y tipo si fuera necesario, y hace clic en \IUbutton{Guardar Fuente}.
	\UCpaso[\UCsist] Verifica que la URL no esté ya registrada \ErrorRef{CU-02}{E1}{URL duplicada}.
	\UCpaso[\UCsist] Crea el registro en la tabla \texttt{news\_sources} con \texttt{is\_active = True} y \texttt{last\_status = 'pending'}.
	\UCpaso[\UCsist] Muestra mensaje de confirmación y actualiza la tabla de fuentes en pantalla.
\end{UCtrayectoria}

%--------------------------------------
\begin{UCtrayectoriaA}{CU-02}{A}{La URL ingresada ya está registrada en el catálogo}
	\UCpaso[\UCsist] Detecta el registro duplicado en la consulta de unicidad.
	\UCpaso[\UCsist] Retorna HTTP 409 con el mensaje ``Esta URL ya existe como fuente''.
	\UCpaso[\UCActor] Lee la notificación de error y puede modificar la URL o desactivar la entrada existente.
	\UCpaso[] El Caso de Uso termina sin crear un nuevo registro.
\end{UCtrayectoriaA}

%--------------------------------------
\begin{UCtrayectoriaA}{CU-02}{B}{El administrador intenta eliminar una fuente predeterminada del sistema}
	\UCpaso[\UCsist] Detecta que el atributo \texttt{is\_predefined = True} está activo para el registro.
	\UCpaso[\UCsist] Retorna HTTP 403 con el mensaje ``No se puede eliminar una fuente predeterminada. Puedes desactivarla en su lugar.''.
	\UCpaso[\UCActor] Opta por usar el interruptor \IUbutton{Desactivar} en lugar de eliminar.
	\UCpaso[] El Caso de Uso termina sin modificar el registro.
\end{UCtrayectoriaA}

%--------------------------------------
\begin{UCtrayectoriaA}{CU-02}{C}{La URL ingresada no tiene esquema HTTP válido}
	\UCpaso[\UCsist] El validador del backend detecta que la cadena no inicia con \texttt{http://} ni \texttt{https://}.
	\UCpaso[\UCsist] Retorna HTTP 400 con mensaje ``La URL debe comenzar con http:// o https://''.
	\UCpaso[\UCActor] Corrige la URL en el formulario y reintenta.
	\UCpaso[] El Caso de Uso continúa en el paso 7 de la trayectoria principal.
\end{UCtrayectoriaA}
